% ch2 B.3 自由电子气中的交换与关联 (原书页 53–60 / PDF p068–p075)

\subsection{自由电子气中的交换与关联}

写到这里，各位也许会有一点惊讶：我们竟然已经非常接近于算出真实的结合能了。你完全有理由发问：金属电子自身的 Coulomb 相互作用又怎么样了呢？表面看来，我们把这些完全忽略了，因为在计算波函数时我们只用了离子实的场。

实际上，价电子之间的相互作用已在很大程度上以一种非常巧妙而不动声色的方式考虑了进来。Wigner 所做的，实质上是把价电子处于不同原子上时的相互作用包括了进来，因为他让近邻的元胞呈电中性，而不像\emph{只}有离子实存在时那样带 $+1$ 电荷；但他有意略去了同一元胞内电子之间的相互作用。对于交换（exchange）与关联（correlation）效应而言，这是一个十分合理的近似。先说交换——我们知道，不相容原理会在一个给定电子周围的自旋平行电子分布中挖出一个“空穴”。Coulomb 排斥则倾向于把自旋反平行的电子也从给定电子的邻域排开，于是净效果是：从该电荷的邻域中大约移走了整整一个电子的电荷。这个空穴的半径为 $r_s$ 的量级，因此在计算波函数时，不妨当作每个元胞内在同一时刻只允许存在一个电子来处理。结果是：无须求解任何非常复杂的波动方程，我们所处理的波函数就已经算到了实际上与 Hartree-Fock 相当的精度。在多价金属中这种空穴相对较小、因而这一近似不再适用——这一事实也许是在此类金属中结合能计算困难的另一个原因。

尽管如此，为了得到真正精确的最终答案，关键还是要回到这些波函数，并尝试对“多体”修正作出更完备的估计。基于上述理由，我们预期所有这些修正之和接近于零，但各个单项却会相当大。

使 Wigner 和 Seitz (28, 29) 能够对这些项作出良好估计的，是一个幸运的事实：碱金属中的电子，除了那个使总能量整体下降的常数“边界修正”之外，其行为非常像一团真正自由的电子气体。如前所述，在 Na 中，在元胞体积的 90\% 以上的区域内，波函数实际上就是 $e^{i\mathbf{k}\cdot\mathbf{r}}$。这使得我们可以使用自由电子气自相互作用能的理论。而这一理论由于气体必须在空间中“均匀”这一事实而大为简化——也就是说，它在空间每一点处的性质必须相同。

人们最先施加的修正是 Hartree 修正，即所有电子的平均场。在自由电子的情形下，它是整个均匀电荷分布的自相互作用能，并按 $v^{2/3}$ 发散；当然，这一发散必定会被来自离子电荷的一个等值反号的项所抵消，正是后者保证了电中性。研究自由电子气时，通常的做法是完全忽略这一项：人们用一团均匀的正电荷“果冻”（jelly）来代替离子实，它恰好抵消自由电子的平均电荷。

而在真实的碱金属中，情况则不同：确实存在这样一项，而且事实上相对于内聚能它还是大的。我们此前已把与中心电荷球以外一切的相互作用包括了进来，所以还必须把一个电荷为 $e$、半径为 $R_{\mathrm{cell}} = r_s a_0$ 的球的 Coulomb 自能包括进来，它等于 $3e^2/5R_c = 6/5r_s$ ryd。如你所见，在 Li 中这超过了 4 伏特，足以把全部符合都破坏掉。

为了与这一项相抗衡，我们必须计算那两种倾向于使电子彼此分开的效应：交换本身，以及关联。后一种能量的定义相当人为：它是最优的 Hartree-Fock 解与正确的解之间的能量差。我们不打算在这里计算关联能（correlation energy），但交换效应我们能够处理；而且，由于其数学与思想日后对我们有用，我们也就来做这件事。

注意到在自由电子气的情形下 Hartree 项只是一个平庸的常数，包含交换的 H-F 方程即为
\begin{equation*}
-\frac{\hbar^2}{2m}\,\nabla^2\phi_{k\sigma} - \sum_{k'\ \mathrm{occ}} \int \mathrm{d}\mathbf{r}'\ \phi^*_{k'\sigma}(\mathbf{r}')\,\phi_{k\sigma}(\mathbf{r}')\,\frac{1}{|\mathbf{r}-\mathbf{r}'|}\,\phi_{k'\sigma}(\mathbf{r}) = E_k\,\phi_{k\sigma}(\mathbf{r})
\end{equation*}
由于系统作为一个整体仍具有平移对称性，通过假设
\begin{equation*}
\phi_k = \frac{1}{\sqrt{V}}\,e^{i\mathbf{k}\cdot\mathbf{r}}
\end{equation*}
便可以得到一个自洽解——即便它不一定就是最低的解。于是交换项变为
\begin{equation*}
A\phi_k(\mathbf{r}) = \int \mathrm{d}\mathbf{r}'\ A(\mathbf{r},\mathbf{r}')\,\phi_k(\mathbf{r}')
\end{equation*}
\begin{align*}
A(\mathbf{r},\mathbf{r}') &= \sum_{k'}\frac{e^2}{|\mathbf{r}-\mathbf{r}'|}\,\phi^*_{k'}(\mathbf{r}')\,\phi_{k'}(\mathbf{r})\\
&= \frac{e^2}{V}\sum_{k'}\frac{e^{i\mathbf{k}'\cdot(\mathbf{r}-\mathbf{r}')}}{|\mathbf{r}-\mathbf{r}'|} = A(\mathbf{r}-\mathbf{r}')
\end{align*}
这确实具有平移对称性（不依赖于 $\mathbf{r}+\mathbf{r}'$），因此方程的解的确可以是平面波。

我们实际上能够把 $A(\mathbf{r}-\mathbf{r}')$ 的形式计算出来；它常被称作交换空穴（exchange hole）的势。这一推导的大部分应归功于 Dirac (30)。我们可以把它看作围绕 $\mathbf{r}'$ 处电子的电荷分布中一个“空穴”所产生的势：
\begin{align*}
\rho(\mathbf{r}-\mathbf{r}') &= \frac{e}{V}\sum_{k\ \mathrm{occ}} e^{i\mathbf{k}\cdot(\mathbf{r}-\mathbf{r}')}\\
&= \frac{e}{V}\left(\frac{n}{2}V\right) \frac{\displaystyle \int_0^{k_F} k^2\,\mathrm{d}k \int_{-1}^{1}\mathrm{d}x\ e^{ik|\mathbf{r}-\mathbf{r}'|x}}{\displaystyle 2\int_0^{k_F} k^2\,\mathrm{d}k}
\end{align*}
利用
\begin{equation*}
\frac{n}{2} = \frac{4\pi}{3}\,\frac{k_F^3}{8\pi^3} = \frac{k_F^3}{6\pi^2}
\end{equation*}
并令 $R = |\mathbf{r}-\mathbf{r}'|$，它就是
% 校注：原书此处两处分母均印作 2\pi^2 R；按本节自洽性（下文 A(R) =
% e^2 k_F^4/2\pi^2[...]，以及精确核 A = (e^2/V)\sum e^{ik·R}/R），
% 分母应为 2\pi^2 R^2，此处按本意转写。
\begin{equation*}
= -\frac{e^2}{2\pi^2 R^2}\,\frac{\mathrm{d}}{\mathrm{d}R}\left\{\int_0^{k_F} \cos kR\ \mathrm{d}k\right\} = -\frac{e^2}{2\pi^2 R^2}\,\frac{\mathrm{d}}{\mathrm{d}R}\left(\frac{\sin k_F R}{R}\right)
\end{equation*}
\begin{align*}
&\left(\text{因为当 } k \to 0 \text{ 时 } \frac{\sin kR}{R} \to 0\right)\\
% 校注：原书前置因子印作 e k_F/2\pi（放大核查确无上标）；但由 \rho(0)=en/2
% （图 12 中水平线标记 n/2 = e k_F^3/6\pi^2，曲线自零出发）以及正文陈述
% A = e\rho/R（与下式 A(R) 精确一致），前置因子应为 e k_F^3/2\pi^2，
% 此处按本意转写。
\rho(R) &= \frac{e k_F^3}{2\pi^2}\left[-\frac{\cos k_F R}{(k_F R)^2} + \frac{\sin k_F R}{(k_F R)^3}\right]
\end{align*}
图 12 给出了交换空穴的图像。（注意在 $R = 0$ 处诸奇性相互抵消。）图中画出的是：一个位于零点的电子在 $R$ 处所看到的自旋平行电子的总有效密度。

%FIG{12}: 图 12　交换空穴：位于原点的电子周围自旋平行电子的总有效密度 $n/2-\rho$ 随 $k_F R$ 的变化，水平线标记 $n/2 = ek_F^3/6\pi^2$（原书图注仅作 “Figure 12”）

相应的核——交换势 $A(\mathbf{r}-\mathbf{r}')$——当然就是 $e\rho$ 除以 $R = \mathbf{r}-\mathbf{r}'$：
\begin{equation*}
A(R) = \frac{e^2 k_F^4}{2\pi^2}\left(\frac{\sin k_F R}{(k_F R)^4} - \frac{\cos k_F R}{(k_F R)^3}\right)
\end{equation*}

这些函数带着它们在无穷远处特有的甚长程振荡行为，在自由电子气理论中极其常见也极其有用。$A$ 类似于自由电子气的“Green 函数”，而熟悉自旋共振的读者会在 $A(\mathbf{r})$ 中认出 Ruderman-Kittel 函数。这种长程的振荡行为，是对费米分布的积分中存在尖锐上限的结果，而且只有在存在自由费米面时才会出现。这两个函数包含的是 $2k_F R$ 而非 $k_F R$，因为它们来自一个稍有不同的积分，但振荡的缘由是相同的。

一个给定平面波态 $k$ 的能量中的交换贡献，很容易由
\begin{equation*}
A\phi_k(\mathbf{r}) = \int A(\mathbf{r}-\mathbf{r}')\,\phi_k(\mathbf{r}')\ \mathrm{d}\mathbf{r}'
\end{equation*}
导出。根据对 Fourier 变换作折叠的卷积定理（convolution theorem），它就是
\begin{equation*}
A\phi_k = A_k\phi_k
\end{equation*}
其中
\begin{equation*}
A_k = \int \mathrm{d}\mathbf{R}\ e^{i\mathbf{k}\cdot\mathbf{R}}\,A(R)
\end{equation*}
这个 Fourier 变换可以直接计算，也可以按下面的办法来做：$A(\mathbf{r}) = \rho(\mathbf{r})/r$，于是再次利用卷积定理，
% 校注：原书此行被积函数印作 \rho(k)；按卷积定理并对照下一行应为 \rho(k')，
% 此处按本意转写。
\begin{align*}
A_k &= \frac{e}{(2\pi)^3}\int \mathrm{d}\mathbf{k}'\ \rho(k')\left(\frac{1}{r}\right)_{k-k'}\\
&= \frac{4\pi e}{(2\pi)^3}\int \mathrm{d}\mathbf{k}'\ \frac{\rho(k')}{|k-k'|^2}
\end{align*}
但 $\rho(k')$ 在 $k'$ 被占据即 $k' < k_F$ 时就是 $e/V$，否则为 0；于是
% 校注：原书左端前置因子印作 4\pi e^2/((2\pi)^3 V)，与其右端 e^2/2\pi^2
% 相差一个 (2\pi)^3 因子，等式不成立；使等式成立的前置因子为 4\pi e^2/V
% （右端 e^2/2\pi^2 \int 为标准结果，且与下文 0.612/r_s 一致），按本意转写。
\begin{equation*}
A_k = \frac{4\pi e^2}{V}\sum_{k'\ \mathrm{occ.}}\frac{1}{|k-k'|^2} = \frac{e^2}{2\pi^2}\int\frac{\mathrm{d}^3k'}{|k-k'|^2}
\end{equation*}
据我所知，这个积分只能用直接的办法费力地算出；这项工作由 Dirac 于 1930 年完成。结果是
\begin{equation*}
A_k = \frac{0.612}{r_s}\left[2 + \frac{k_F^2 - k^2}{k\,k_F}\,\ln\left|\frac{k+k_F}{k-k_F}\right|\right]
\end{equation*}
这个积分是一个非常著名而且有用的函数。我们可以把方括号中的量勾画如下：它在 $k = 0$ 处的值为 4；分母中的奇性被 $\ln|k+k_F/k-k_F| \to 2k/k_F$ 所抵消。在 $k = k_F$ 处存在第二个奇性，但它取 $0\ln 0 = 0$ 的形式，故 $A = 2$，只是该点的斜率为无穷大（图 13）。

我们可以把这条曲线看作费米分布函数的一幅被抹糊了的复本，其无穷大的斜率正对应于费米面处的陡然下降。其原因在于：当 $k < k_F$ 时，分母在 $k' = k$ 处趋于零；而当 $k > k_F$ 时则不然。这个零点虽是可积的，但也仅仅是勉强可积而已。这反过来又反映出 Coulomb 势的长程性质：$1/r$ 的 Fourier 变换是 $2\pi/k^2$。\footnote{无穷大的斜率可以直接证明：取
\begin{align*}
\nabla_k A_k &= \frac{4\pi e^2}{V\,\delta k}\left[\sum_{k\ \mathrm{occ.}}\frac{1}{|k-k'|^2} - \sum_{(k+\delta k)\ \mathrm{occ.}}\frac{1}{|k-k'|^2}\right]\\
&= -\frac{4\pi e^2}{V}\sum_{\text{费米球面}}\frac{\cos\theta}{|k-k'|^2}
\end{align*}
面元按 $k\,\mathrm{d}k$ 变化，于是得到对数发散 $\int k\,\mathrm{d}k/k^2$。}

%FIG{13}: 图 13　$A_k$ 表达式中方括号内的函数随 $k$ 的变化：$k=0$ 处值为 4，$k=k_F$ 处斜率无穷大，横轴标出 $k$、$k_F$、$2k_F$（原书图注仅作 “Figure 13”）

由 $A(k)$ 我们可以计算两件事。第一件是平均交换能（exchange energy），它将是 $(1/2N)\sum_k A(k)$。我们看到 $A$ 在 $0.612/r_s$ 的（2 到 4）倍之间变化；平均因子约为 3：$E_{\mathrm{exch}} = -0.916/r_s$，它抵消掉了 $1.2/r_s$ 的 Coulomb 自能中很大的一部分。尽管如此，在 Na 中 $0.25/r_s$ ryd $\cong 0.85$ ev 的差额，已足以移去内聚能的相当大一部分。

剩下的修正是“关联能”，即由于即使是自旋反平行的电子也倾向于彼此保持距离而必须作出的那项修正。关于这一修正项有着浩繁的文献；只需说明这一点：Wigner 以及后来的 Pines (31)，在一类对 $r_s \lesssim 1$ 肯定良好的表达式与另一类对 $r_s \gtrsim 10$ 同样良好的表达式之间，于相当宽的范围内作了内插——尽管对于 $r_s = 4$ 而言，这两类表达式都谈不上有合理的希望能够接近——由此得到如下估计。鉴于它在已知范围内变化不快，而且各种估计彼此相符，我们可以猜测它在 $\pm 20\%$ 的范围内有效：$E_{\mathrm{corr}} = 0.88/(r_s + 7.8)$ ryd。［许多作者所用的并非此式，而是 Wigner 更早的一个表达式；他于 1938 年在一条脚注中对其作了修正 (29)。］用这个表达式我们得到表 4 所给出的结果。

\begin{center}
表 4

\begin{tabular}{lccccc}
\toprule
元素 & 未修正 $E_{\mathrm{coh}}$，ev & $\Delta E_{\mathrm{HF}}$，ev & $\Delta E_{\mathrm{corr}}$，ev & 总计，ev & 实验，ev\\
\midrule
Li & $-2.1$ & $1.19$ & $-1.09$ & $-2.0$ & $-1.65$\\
Na & $-1.15$ & $0.96$ & $-1.02$ & $-1.2$ & $-1.2$\\
Rb & $-0.42$ & $0.73$ & $-0.92$ & $-0.61$ & $-0.85$\\
\bottomrule
\end{tabular}
\end{center}

我必须提醒各位：我在这里给出的这些表并不是权威的，也不代表 Q.D.M. 所已取得的最精确或最好的结果。至于更详细的结果，我请各位去查阅原始文献，特别是 Ham 的文章以及 Brooks 收在 Varenna 夏季学校报告中的文章。正如这后一篇文章所指出的，Ham 所取得的良好符合会因各种进一步的修正而打折扣。能够可靠引出的结论也许只有如下这些：Wigner 对自由电子关联能的估计离正确值并不远；而且我们确实知道碱金属中内聚的真正物理来源。与关联能中 $\pm 0.2$ ev 的误差同量级的修正，可能是费米能中的其他项，而最重要的则是由于波函数并非自由电子波函数而引起的对交换与关联的修正。稍后我们会看到，仅仅在这些效应的公式中塞进一个 $m^*$ 是不合法的，而这些修正的真正本质目前尚不清楚。［不过，可参见不久后将要讨论的 Phillips 的工作 (25)。］这些修正很容易达到关联能的 20\% 的量级。

Q.D.M. 的数值结果中有一个方面引人注目且非常令人鼓舞：在这些金属的压缩率上得到的符合，几乎与在能量上得到的符合同样好。早年的许多计算以及大多数用其他方法所作的计算，给出的压缩率都非常糟糕。毕竟，压缩率是结合能曲线的二阶差分。Q.D.M. 之所以如此之好，其原因可能在于：对所有各项随密度的依赖形式我们确实了解得相当清楚——$E_0$ 项是因为它可以算得非常精确，其余各项则是因为它们由一般性的理论考虑所决定。动能按 $1/r_s^2$ 变化，而交换能、Coulomb 能与关联能则大致按 $1/r_s$ 变化。基于这一想法，用 Bardeen 的半经验状态方程 (27) 常常能够对 pV 曲线得到非常好的拟合：
\begin{equation*}
E = \frac{A}{r_s^3} + \frac{B}{r_s^2} - \frac{C}{r_s}
\end{equation*}

在接着讨论 O.P.W. 方法之前，我现在想先回到交换能上来略作讨论。$A(k)$ 的另一个重要方面是：它代表着对动量为 $k$ 的电子的自能（self-energy）$E_k$ 的一个依赖于 $k$ 的修正：$E_k = \hbar^2 k^2/2m^* - A(k)$。这是诸如交换这类非局域（nonlocal）势的特征行为。对于由各种多体效应所产生的有效势来说，非局域性是常态。这个势是 $A\phi = \int A(\mathbf{r}-\mathbf{r}')\,\phi(\mathbf{r}')\ \mathrm{d}\mathbf{r}'$；由卷积定理，它变换成自能的一个依赖 $k$ 的修正 $A\phi_k = A_k\phi_k$，可以粗略地把它看作一种\emph{有效质量}修正——不过它并非由晶格的周期势所引起，而是由交换空穴随粒子动量的变化所引起。

在采用 Coulomb 相互作用的交换这一特殊情形下，有效质量修正恰好在费米面处发散：由于 $A(k)$ 具有无穷大的负斜率，$\mathrm{d}E_k/\mathrm{d}k\big|_{E_F} = +\infty$ 恰好在费米面处成立。态密度对数式地降到零。假如果真如此，就会与低温比热、自旋顺磁性等方面的许多实验结果相矛盾。Wigner 在原始论文中就已指出：实际上，在真实电子气中，Coulomb 相互作用的长程部分会被近旁其他电子的微小位移所屏蔽（screening），而这些位移同样也牵涉在关联能之中。对这种势屏蔽的一个粗略计算给出有效势 $e^{-k_0 r}/r$，其中 $k_0$ 是 Debye 波数，其量级为 $k_F$。这样一来，$A(k)$ 中的奇性被抹平，有效质量修正不再发散。我们应当强调，“交换空穴”仍然是一个长程的、振荡的函数，在 $k$ 空间的费米面处带有奇性；但它对交换势 $A$ 的影响，则会因相互作用范围的缩短而被抹平。这时 $A(k)$ 对 $k$ 的曲线看上去就像图 14 那样。

%FIG{14}: 图 14　交换势 $A(k)$ 与屏蔽后的 $A_{\mathrm{eff}}$ 随 $k$ 的变化，水平线为 Slater 近似，$k_F$ 处的奇性已被抹平（原书图注仅作 “Figure 14”）

% ---- B.3 收尾（原书 p.61 / PDF p076 顶部，协调者补译） ----
其净结果是有效质量的减小，在金属密度下其数值可能不大（5\% 到 20\%）。

在实际晶体中，自由电子的密度并不是常数，自由电子模型不再适用，$A$ 也不再是 $A(\mathbf{r}-\mathbf{r}')$ 而是 $A(\mathbf{r},\mathbf{r}')$；它既是非定域的，又是随空间变化的。一个常常采用、源自 Slater (32) 的近似，是在这种情形下忽略 $A$ 的非定域性质，试图把它换成最好的定域势 $A(\mathbf{r})\,\delta(\mathbf{r}-\mathbf{r}')$，后者将正比于局域密度的 $1/3$ 次幂。$\delta$ 函数的傅里叶变换是常数，所以这样做的实质，是用一个适当的平均值代替 $A(\mathbf{k})$（见图~14）。这一近似能够成立的判据是：电子密度相对于交换空穴的尺度变化缓慢，而交换空穴的尺度大致为 $r_s$；此时非定域势所取样的区域将具有大致恒定的交换势。这样的近似在碱金属中是不可想象的——在那里 $r_s > r_{\mathrm{core}}$，因而远大于发生变化的长度尺度——但在 $r_s \sim 1$--$2$ 的多价物质中情形可能要好一些。Phillips 对金刚石（$r_s = 1.3$ (33)）就这一问题作过相当透彻的研究，结果表明 Slater 近似在那里并不严重，尽管对能带结构的某些方面仍有 20\% 到 30\% 的修正，尤其是对从最低态 $k=0$、$L=0$ 量起的费米面整体高度的修正。当然，这正是永远可以预期的一类效应：它对应于真实自由电子气中 $A(\mathbf{k})$ 的动量依赖性。\footnote{关于把核物质当作费米气体的最早评论，大概出自 Van Vleck：他指出在这种情形下，交换对有效质量的修正同样可能导致较轻的 $m^*$——这是一个确实发生并且至关重要的效应。}
